\scanpage{78}
% archive foliation: 77
\uline{Relations avec op.} $\otimes$ \uline{et} $\Gamma^p$.

La donnée d'une cat. polynomiale \inserted{$\unsure{\mathcal P|P}$} permet de définir un foncteur triadditif
\[
\unsure{\mathcal P|P}_1^\circ\times \unsure{\mathcal P|P}_1^\circ\times \unsure{\mathcal P|P}_1\longrightarrow\operatorname{Ab}
\]
par
\[
(X,Y;Z)\longmapsto \unsure{\mathcal P|P}_{(1,1)}(X,Y;Z)
\]
(et de même pour plus d'arguments de gauche).
Si, pour $X,Y$ fixés, $Z$ variable, ce foncteur est représentable, l'objet qui le représente est noté $X\otimes Y$. Il dépend fonctoriellt de $X,Y$ dans $\unsure{\mathcal P|P}_1$ (mais non pour $X,Y$ \deleted{\unsure{$\in\operatorname{Ob}\unsure{\mathcal P|P}$}} considérés comme dans $\unsure{\mathcal P|P}$ !). \deleted{\unsure{Il}} On sait que $X\otimes Y$ existe ssi $Y\otimes X$ existe, et \unsure{il y a un} can. isom (noté $s_{X,Y}$), on a
\[
s_{\unsure{X\otimes Y|X,Y}}\circ s_{\unsure{Y\otimes X|Y,X}}=\operatorname{id}_{Y\otimes X}.
\]

\unsure{Sur|Pour} $X,n\in\mathbf N$ fixés, si le foncteur
\[
Y\longmapsto \unsure{\mathcal P|P}_n(X,Y)
\]
sur $\unsure{\mathcal P|P}_1$ est représentable, l'objet qui le représente est noté $\Gamma^n X$; il dépend encore fonctoriellt de $X$ dans $\unsure{\mathcal P|P}_1$. L'objet $\Gamma^0(X)$ existe ssi $\Gamma^0(0)$ existe il est noté $1$.

\unsure{Prouvons} que \deleted{\unsure{$\forall$}} si $1$ existe, alors $\forall X\in\operatorname{Ob}\unsure{\mathcal P|P}$, $1\otimes X\simeq X$ i.e.
\[
\unsure{\mathcal P|P}_{(1,1)}(1,X;Y)\mathrel{\mathop{\simeq}\limits_{\unsure{\text{isom bifonctoriel}|\text{sous bifoncteurs}}}}\unsure{\mathcal P|P}_1(X,Y)\bigl(=\operatorname{Hom}_{\unsure{\mathcal P|P}_1}(X,Y)\bigr)\quad ?
\]
et plus particulièrement $1\otimes1\simeq1$ i.e. $\forall Y$
\[
\unsure{\mathcal P|P}_{1,1}(1,1;Y)\simeq \unsure{\mathcal P|P}_1(1,Y)\qquad ?
\]

